\subsection{Negative displaced contact: the weak root on the linear scale}
\label{sec:direct-negative-displaced}

The preceding estimate loses its normalized constant-coefficient
margin when that coefficient becomes small. If
\(q_2/\sqrt{q_1}\) remains bounded, the approaching positive boundary
root gives the weak-root scale. \Cref{prop:weak-root-passage} reduces
the count to an estimate for the negative branch beyond the weak-root
radius. We verify this estimate on the linear scale.

\begin{proposition}[Negative displaced contact]
\label{prop:direct-negative-displaced}
For every fixed \(A_0<\infty\), there are fixed
\(a_{\rm w},\kappa_{\rm w}>0\) such that the return displacement has
at most two distinct zeros on
\[
 q_1=\rho_1^2>0,\quad q_2=\bar c\rho_1,\quad |\bar c|\le A_0,\quad
 q_0=-a\rho_1^3,\quad 0<a\le a_{\rm w},\quad
 0<\kappa=\rho_1/a\le\kappa_{\rm w} ,
\]
for all sufficiently small \(\rho_1\) and original parameters.
\end{proposition}

\begin{proof}
To apply \cref{prop:weak-root-passage}, set
\(\ell=\rho_1a=|q_0|/q_1\). Then
\eqref{eq:weak-root-scales} holds with \(\hat r=1\) and
\(\epsilon_2=-\bar c a\), \(|\epsilon_2|\le A_0a\), and the common
logarithmic height coordinate is \(T=-\rho_1\log h=-\kappa a\log h\).
At the radius \(\ell_d=\rho_1d\) the negative drift is
\[
 K_d(x)=a_d+x-\bar c dx^2+d^2x^3R_q(-\ell_dx),
\]
so \eqref{eq:weak-root-negative-drift} holds for a small fixed
\(d_{\max}=d_0\) depending on \(A_0\).

It remains to verify \eqref{eq:weak-root-outer} for the negative graphs
with \(z<Y_b|\eta|^2\) at \(|\eta|=\rho_1d_0\). Return to the linear
scale \(\rho_1\), with \(|\eta|=\rho_1X\). Since \(\ell/a=\rho_1\),
the negative drift on this scale is \(K_1\) from
\cref{cor:weak-root-negative-first-root}:
\[
 K_-^{\rho_1}(X)=a+X-\bar c X^2+X^3R_q(-\rho_1 X).
\]
Its first positive root \(\beta\) stays in a fixed compact interval.
At \(a=\rho_1=0\), it is the positive root of
\(1-\bar c X-X^2=0\); the bound \(|\bar c|\le A_0\) gives uniform
separation and a simple-root margin. On \([d_0/2,\beta]\),
\((K_-^{\rho_1})'(X)/X\ge-C_{A_0,d_0}\).

Apply \cref{cor:weak-root-negative-first-root} with
\(d_{\max}=d_0\) and with \(\beta^*\) the right end of that compact
interval. It gives \(\mathfrak H_-\ge-C_{A_0,d_0}-1\) and
\(|I_{-,T}|=o(1)\) for these graphs, including the limiting negative
root. Thus \eqref{eq:weak-root-outer} holds.
\Cref{prop:weak-root-passage} now gives at most two distinct zeros
of \(\Delta_p\), with \(a_{\rm w}\) and \(\kappa_{\rm w}\) equal to
the constants \(a_\circ\) and \(\kappa_\circ\) of that proposition.
\end{proof}
